\section{Related Work}

{\textbf{Self-Supervised Learning:}} Self-Supervised Learning is aimed at learning rich representations of high dimensional unlabeled data to be useful for a wide variety of tasks. The fields of natural language processing and computer vision have seen dramatic advances in self-supervised methods such as BERT \cite{devlin2018bert}, CPC, MoCo, SimCLR \cite{henaff2019data, he2019momentum,chen2020simclr}.

{\textbf{Contrastive Learning:}} Contrastive Learning is a framework to learn representations that obey similarity constraints in a dataset typically organized by similar and dissimilar pairs. This is often best understood as performing a dictionary lookup task wherein the positive and negatives represent a set of keys with respect to a query (or an anchor). A simple instantiation of contrastive learning is Instance Discrimination \cite{wu2018unsupervised} wherein a query and key are positive pairs if they are data-augmentations of the same instance (example, image) and negative otherwise. A key challenge in contrastive learning is the choice of negatives which can decide the quality of the underlying representations learned. The loss functions used to contrast could be among several choices such as InfoNCE \cite{oord2018representation}, Triplet \cite{wang2015unsupervised}, Siamese \cite{chopra2005learning} and so forth.

{\textbf{Self-Supervised Learning for RL}}: Auxiliary tasks such as predicting the future conditioned on the past observation(s) and action(s) \cite{jaderberg2016reinforcement, shelhamer2016loss, oord2018representation, schmidhuber1990a} are a few representative examples of using auxiliary tasks to improve the sample-efficiency of model-free RL algorithms. The future prediction is either done in a pixel space \cite{jaderberg2016reinforcement} or latent space \cite{oord2018representation}. The sample-efficiency gains from reconstruction-based auxiliary losses have been benchmarked in \citet{jaderberg2016reinforcement, higgins2017darla, yarats2019improving}. Contrastive learning has been used to extract reward signals in the latent space \cite{sermanet2018time, dwibedi2018learning, warde2018unsupervised}; and study representation learning on Atari games by \citet{anand2019unsupervised}.

{\textbf{World Models for sample-efficiency:}} While joint learning of an auxiliary unsupervised task with model-free RL is one way to improve the sample-efficiency of agents, there has also been another line of research that has tried to learn world models of the environment and use them to sample rollouts and plan. An early instantiation of the generic principle was put forth by \citet{sutton1990integrated} in Dyna where fictitious samples rolled out from a learned world model are used in addition to the agent's experience for sample-efficient learning. Planning through a learned world model~\cite{srinivas2018universal} is another way to improve sample-efficiency. While \citet{jaderberg2016reinforcement, oord2018representation, lee2019stochastic} also learn pixel and latent space forward models, the models are learned to shape the latent representations, and there is no explicit Dyna or planning.
Planning through learned world models has been successfully demonstrated in \citet{ha2018world, hafner2018learning, hafner2019dream}. \citet{kaiser2019model} introduce SimPLe which implements Dyna with expressive deep neural networks for the world model for sample-efficiency on Atari games.

{\textbf{Sample-efficient RL for image-based control:}} CURL encompasses the areas of self-supervision, contrastive learning and using auxiliary tasks for sample-efficient RL. We benchmark for sample-efficiency on the DMControl suite \citep{tassa2018deepmind} and Atari Games benchmarks \cite{bellemare2013arcade}. The DMControl suite has been used widely by \citet{yarats2019improving}, \citet{hafner2018learning}, \citet{hafner2019dream} and \citet{lee2019stochastic} for benchmarking sample-efficiency for image based continuous control methods. As for Atari, \citet{kaiser2019model} propose to use the 100k interaction steps benchmark for sample-efficiency which has been adopted in \citet{kielak2020rainbow,  van2019use}. The Rainbow DQN \cite{hessel2017rainbow} was originally proposed for maximum sample-efficiency on the Atari benchmark and in recent times has been adapted to a version known as Data-Efficient Rainbow \cite{van2019use} with competitive performance to SimPLe without learning world models. We benchmark extensively against both model-based and model-free algorithms in our experiments. For the DMControl experiments, we compare our method to Dreamer, PlaNet, SLAC, SAC+AE whereas for Atari experiments we compare to SimPLe, Rainbow, and OverTrained Rainbow (OTRainbow) and Efficient Rainbow (Eff. Rainbow).
